\section{Rust and MIR}
\label{sec:rust-and-mir}

This section briefly situates Rust relative to Java~\cite{gustedt_2002}, describes the features of Rust relevant to treewidth, and introduces MIR, the compiler's internal CFG representation on which our measurements are based.
Rust is a systems programming language designed to combine low-level control over memory with strong compile-time safety guarantees.
Unlike Java, Rust does not rely on a garbage collector; instead, memory safety is enforced statically at compile time by the \emph{borrow checker}~\cite[Ch.~4.2]{rust_book}.

\subsection{Ownership and Borrowing}
\label{subsec:ownership-borrowing}

Every value in Rust has a unique \emph{owner}~\cite[Ch.~4.1]{rust_book}.
Ownership can be transferred by assignment or by passing a value to a function.
After such a transfer, the original variable may no longer be used, and when the owner goes out of scope, the value is automatically deallocated.
This prevents multiple independent variables from being responsible for the same piece of memory.

Programs sometimes need temporary access to a value without transferring ownership.
Rust supports this through \emph{borrowing}~\cite[Ch.~4.2]{rust_book}: a borrow is a reference to a value owned elsewhere.
There are two kinds: \emph{shared borrows} (\texttt{\&T}), which allow read-only access, and \emph{mutable borrows} (\texttt{\&mut T}), which allow modification.
Rust enforces strict aliasing rules: any number of shared borrows may coexist, or exactly one mutable borrow may exist, but not both at the same time.
A borrow is only valid for some region of the program, called its \emph{lifetime}~\cite[Ch.~10.3]{rust_book}.
References may not outlive the values they refer to, preventing dangling pointers statically.
Together, these aliasing and lifetime constraints are part of Rust's static semantics, enforced at compile time by the borrow checker.
These checks are \emph{flow-sensitive}: whether a borrow is active depends on the control-flow path taken through the program.
The borrow checker therefore operates directly on the CFG, and ownership constraints can influence its structure: values that go out of scope on some paths but not others require explicit drop edges in the compiled control-flow graph.

\subsection{Rust Control Flow}
\label{subsec:rust-control-flow}

The structure of Rust's control flow is the other main factor shaping CFG complexity.
Rust uses predominantly \emph{structured} control flow: execution is expressed through conditional expressions, pattern matching, and loops rather than arbitrary jumps~\cite[Ch. 8.2.13]{rust_reference}.
This is one of the main reasons CFGs of real programs tend to remain relatively simple, as prior work on goto-free language fragments has shown~\cite{thorup_1998}.

Certain language features can make the compiled control-flow graph more complex than the source structure suggests.
Labeled \texttt{break} and \texttt{continue} can jump to an enclosing loop, creating edges in the CFG that bypass intermediate blocks.
Gustedt et al.\ showed that such constructs are similarly harmful to \texttt{goto} and can lead to arbitrarily high treewidth~\cite{gustedt_2002}.
Similarly, async functions compile to state machines whose branches and cycles are not visible at the source level.
This motivates measuring treewidth empirically on the compiler's internal representation (see Section~\ref{subsec:mir}) rather than reasoning from source-level nesting alone.

\subsection{Mid-level Intermediate Representation}
\label{subsec:mir}

The Rust compiler lowers functions into the \emph{Mid-level Intermediate Representation} (MIR) after parsing and earlier compilation stages~\cite[Ch.~44]{rustc_dev_guide}.
MIR is a control-flow graph: it consists of basic blocks containing simple typed statements, together with explicit control-flow transfers between those blocks.
A basic block is a sequence of statements followed by a \emph{terminator}.
Terminators represent control transfers with potentially multiple successors, covering branching, function calls, drops, and returns.
Once lowered to MIR, execution order is no longer implicit in the source structure but is represented explicitly as a graph.

MIR provides a direct basis for studying the structure of real Rust programs.
It offers a CFG representation suitable for treewidth analysis, and coincides with the level at which the compiler performs its own flow-sensitive analyses~\cite[Ch.~44]{rustc_dev_guide}.
We therefore extract and analyse CFGs at the MIR level, where the compiler's own structural reasoning takes place.
The next section covers this in detail.
